% !TEX program = xelatex
\documentclass[a4paper,12pt]{report}

% XeLaTeX 與繁體中文排版
\usepackage{fontspec}
\usepackage{xeCJK}
\setCJKmainfont[AutoFakeBold=2,AutoFakeSlant=0.2]{標楷體}
\setCJKsansfont{Microsoft JhengHei}
\setCJKmonofont{SimSun}
\setmainfont{Times New Roman}
\XeTeXlinebreaklocale "zh"
\XeTeXlinebreakskip = 0pt plus 1pt

% 版面、表格、程式碼與連結
\usepackage[margin=2.4cm,headheight=15pt]{geometry}
\usepackage{array}
\usepackage{listings}
\usepackage{xcolor}
\usepackage{hyperref}
\usepackage{fancyhdr}
\usepackage{enumitem}
\hypersetup{colorlinks=true,linkcolor=blue!55!black,urlcolor=blue!65!black,citecolor=blue!55!black,pdftitle={LINE Messaging API 註冊申請與實作教學},pdfauthor={簡宏洲}}
\setlength{\parindent}{2em}
\setlength{\parskip}{0.35em}
\setlength{\emergencystretch}{2em}
\setlist{itemsep=0.2em,topsep=0.3em}
\pagestyle{fancy}
\fancyhf{}
\fancyhead[L]{LINE Messaging API 教學文件}
\fancyhead[R]{學號 413631655}
\fancyfoot[C]{\thepage}
\renewcommand{\headrulewidth}{0.4pt}

% 中文標題與章節編號
\renewcommand{\contentsname}{章節目錄}
\renewcommand{\listfigurename}{圖目錄}
\renewcommand{\listtablename}{表目錄}
\renewcommand{\bibname}{參考資料}
\renewcommand{\chaptername}{第}
\renewcommand{\thechapter}{\arabic{chapter}章}
\renewcommand{\thesection}{\arabic{chapter}.\arabic{section}}
\renewcommand{\thesubsection}{\thesection.\arabic{subsection}}
\renewcommand{\figurename}{圖}
\renewcommand{\tablename}{表}
\renewcommand{\lstlistingname}{程式碼}
\renewcommand{\lstlistlistingname}{程式目錄}

% listings：左側行號、換行與清晰的程式碼樣式
\definecolor{codeblue}{RGB}{30,72,125}
\definecolor{codegreen}{RGB}{36,112,75}
\definecolor{codegray}{RGB}{100,100,100}
\lstdefinestyle{tutorial}{
  basicstyle=\ttfamily\small,
  keywordstyle=\color{codeblue}\bfseries,
  commentstyle=\color{codegreen},
  stringstyle=\color{red!65!black},
  numbers=left,
  numberstyle=\tiny\color{codegray},
  numbersep=8pt,
  stepnumber=1,
  showstringspaces=false,
  breaklines=true,
  breakatwhitespace=false,
  columns=fullflexible,
  keepspaces=true,
  frame=single,
  tabsize=4,
  captionpos=b,
  literate={\#}{{\#}}1
}
\lstset{style=tutorial}
\renewcommand{\thefigure}{\arabic{chapter}.\arabic{figure}}
\renewcommand{\thetable}{\arabic{chapter}.\arabic{table}}
\AtBeginDocument{\renewcommand{\thelstlisting}{\arabic{chapter}.\arabic{lstlisting}}}

\title{LINE Messaging API 註冊申請與實作教學\\[0.6em]\large 從官方帳號建立到 Webhook 機器人}
\author{簡宏洲\\學號：413631655}
\date{\today}

\begin{document}
\maketitle
\thispagestyle{empty}
\begin{center}
\vfill
\large 技術文件・繁體中文版
\vfill
\end{center}
\clearpage
\pagenumbering{roman}
\tableofcontents
\clearpage
\listoffigures
\clearpage
\listoftables
\clearpage
\lstlistoflistings
\clearpage
\pagenumbering{arabic}

\chapter{LINE Messaging API 教學文件製作}

\section{題目說明}

本文件以繁體中文介紹 LINE Messaging API 的註冊申請、頻道設定、Webhook 串接與訊息收發，並以 Python、Flask 及 LINE 官方 Python SDK v3 建立可執行的文字回聲機器人。讀者完成操作後，能建立 LINE 官方帳號，啟用 Messaging API，讓公開 HTTPS 網址接收 LINE 平台事件，驗證請求簽章，再使用事件中的回覆權杖回覆訊息。

本文是入門教學而非特定帳號的代辦申請。不同地區的官方帳號功能、方案、費用與介面可能變動，請以 LINE Developers 官方文件與帳號管理畫面顯示的最新資訊為準。所有憑證均使用環境變數示範；請勿把真實密鑰放入文件、程式碼或公開版本庫。

\section{學習目標與流程總覽}

閱讀本章後，預期能夠：
\begin{enumerate}
  \item 說明 LINE 官方帳號、Provider、Messaging API Channel 與 Bot Server 的角色。
  \item 按照現行流程建立 Business ID 與官方帳號，並啟用 Messaging API。
  \item 區分 Channel Secret 與 Channel Access Token，安全設定兩項憑證。
  \item 以 HTTPS Webhook 接收文字事件，驗證簽章並使用 reply token 回覆。
  \item 在本機啟動範例，透過 ngrok 暫時公開端點並完成測試。
\end{enumerate}

\begin{figure}[htbp]
\centering
\renewcommand{\arraystretch}{1.35}
\begin{tabular}{>{\centering\arraybackslash}p{0.20\textwidth}c>{\centering\arraybackslash}p{0.20\textwidth}c>{\centering\arraybackslash}p{0.20\textwidth}}
\fbox{\shortstack{使用者\\LINE App}} & $\longrightarrow$ & \fbox{\shortstack{LINE 平台\\Messaging API}} & $\longrightarrow$ & \fbox{\shortstack{Bot Server\\Flask Webhook}} \\
 & & \multicolumn{3}{c}{\small 回覆時由 Bot Server 呼叫 Reply API，經 LINE 平台送回使用者} \\
\end{tabular}
\caption{LINE Messaging API 的基本訊息往返流程}
\label{fig:architecture}
\end{figure}

如圖\ref{fig:architecture}，使用者傳送訊息後，LINE 平台以 HTTPS POST 將 Webhook 事件送至伺服器；伺服器驗證請求來源並解析事件，再向 Messaging API 發出回覆請求。Webhook URL 是事件的接收端點，Access Token 則用於授權伺服器呼叫 LINE API，兩者用途不同。

\chapter{Messaging API 簡介與應用場景}

\section{服務概念與元件}

Messaging API 是 LINE Platform 提供的訊息介面，供開發者建立能與使用者互動的 Bot。Bot 通常綁定一個 LINE Official Account，並在 Developers Console 中以 Messaging API Channel 管理設定。Provider 是管理頻道的開發者或組織容器；使用者、官方帳號與頻道共同構成從聊天介面到後端服務的連線。

\begin{description}
  \item[LINE Official Account] 使用者加入好友及互動的官方帳號，也是訊息服務的入口。
  \item[Provider] 管理 Channel 的個人或組織。建立官方帳號並啟用 Messaging API 時，選擇 Provider 前應確認歸屬，因為此關聯日後不能任意移轉。
  \item[Messaging API Channel] 提供 Channel Secret、Access Token 及 Messaging API 設定的技術通道。
  \item[Bot Server] 開發者維運的 Web 服務，接收事件、執行商業邏輯，再呼叫 LINE API。
  \item[Webhook] LINE 平台對 Bot Server 發出的 HTTP POST 通知，內容包含事件資料及供回覆使用的 reply token。
\end{description}

\section{常見應用情境}

Messaging API 可用於自動客服、訂單與物流通知、預約提醒、活動報名、會員服務、常見問題查詢以及群組內的輔助服務。機器人可回覆文字，也能依 LINE 訊息規格傳送圖片、貼圖、位置、範本或 Flex Message。應用程式應只處理必要資料，清楚告知使用者資料用途，並依所在地區法規與平台政策管理資料留存及推播頻率。

\section{訊息模式與使用時機}

回覆訊息使用事件的 reply token，適合回應剛發生的使用者互動；推播訊息則可依目標對象主動發送，但必須遵守官方帳號的訊息額度與使用者互動規則。多種傳送方式的適用條件不同，正式開發前應查閱官方 API 參考文件與帳號方案。

\begin{table}[htbp]
\centering
\caption{常見訊息傳送方式比較}
\label{tab:message-types}
\begin{tabular}{|p{0.18\textwidth}|p{0.25\textwidth}|p{0.42\textwidth}|}
\hline
\textbf{方式} & \textbf{目標／觸發} & \textbf{常見用途} \\
\hline
Reply 回覆 & 由 Webhook 事件及 reply token 觸發 & 回覆問題、確認操作；回覆權杖有使用時效，不宜延遲處理 \\
\hline
Push 推播 & 指定已知的使用者或聊天目標 & 訂單狀態、預約提醒；通常計入官方帳號可傳送訊息額度 \\
\hline
Broadcast 廣播 & 向符合條件的好友發送 & 重要公告或活動訊息；應克制使用並注意訂閱方案 \\
\hline
\end{tabular}
\end{table}

官方文件指出，Messaging API 的費用與免費訊息數會依國家或地區及官方帳號方案而異。回覆訊息與推播等訊息的計算方式也可能不同；請勿直接將其他地區的價格表套用至台灣帳號，應於官方帳號管理後台確認目前額度與方案。

\chapter{帳號註冊與 Messaging API Channel 建立}

\section{事前準備}

開始前準備一個可登入的 LINE 帳號或電子郵件、官方帳號的名稱及基本資訊、管理者聯絡信箱，以及可安裝 Python 套件的電腦。若代表公司或團隊建立，請先確認帳號管理責任、官方帳號歸屬與 Provider 選擇方式；避免把不相關公司或專案的 Channel 混在個人 Provider 之下。

\section{註冊 Business ID 並建立官方帳號}

\begin{enumerate}
  \item 開啟 LINE Business ID 註冊入口（參考資料\ref{ref:getting-started}），選擇以 LINE 帳號或電子郵件註冊，並完成必要的驗證。
  \item 依頁面指示填寫官方帳號申請資料，例如帳號顯示名稱、業務類別、國家／地區及聯絡資訊。實際欄位依當下介面為準。
  \item 登入 LINE Official Account Manager，確認新建立的官方帳號已出現在帳號清單中，並完成基本個人資料設定。
\end{enumerate}

\section{啟用 Messaging API}

現行官方流程是先建立 LINE Official Account，再從 LINE Official Account Manager 的帳號設定啟用 Messaging API。啟用時系統會建立對應的 Messaging API Channel；不應再依照舊教學嘗試直接從 Developers Console 建立 Messaging API Channel，因為該建立方式已於 2024 年停止。

\begin{enumerate}
  \item 在 Official Account Manager 選取帳號，開啟設定中的 Messaging API 項目並選擇啟用。
  \item 若系統要求開發者資料，依提示填寫姓名與電子郵件，完成 LINE Developers 開發者註冊。
  \item 選擇要管理此帳號的 Provider。若要與既有 LINE Login 或其他 Channel 協作，選擇前應確認它們隸屬於同一個正確 Provider。Provider 指派後不易更換。
  \item 登入 LINE Developers Console，開啟所選 Provider，確認 Messaging API Channel 已建立。進入該 Channel 後，即可檢視基本設定與 Messaging API 設定。
\end{enumerate}

\begin{table}[htbp]
\centering
\caption{帳號與頻道建立流程檢核表}
\label{tab:registration}
\begin{tabular}{|p{0.17\textwidth}|p{0.39\textwidth}|p{0.32\textwidth}|}
\hline
\textbf{階段} & \textbf{操作} & \textbf{確認結果} \\
\hline
Business ID & 註冊或登入 Business ID & 能登入管理介面 \\
\hline
官方帳號 & 填表建立 LINE Official Account & Manager 清單可見帳號 \\
\hline
Messaging API & 從 Manager 設定中啟用 & 系統建立 Messaging API Channel \\
\hline
Provider & 選擇管理該頻道的 Provider & Developers Console 可找到 Channel \\
\hline
\end{tabular}
\end{table}

\chapter{Channel 憑證與安全設定}

\section{Channel Secret 與 Access Token 的差異}

Channel Secret 是 Webhook 簽章驗證所用的共享秘密；LINE 以 Channel Secret 對原始 Webhook 內容計算 HMAC-SHA256 簽章，伺服器再驗證請求標頭中的 \texttt{X-Line-Signature}。Channel Access Token 則是伺服器呼叫 Messaging API 的授權憑證，例如呼叫 Reply API。兩者皆屬機密，但不可互相替代。

\begin{table}[htbp]
\centering
\caption{兩種 Channel 憑證的用途與管理方式}
\label{tab:credentials}
\begin{tabular}{|p{0.19\textwidth}|p{0.31\textwidth}|p{0.37\textwidth}|}
\hline
\textbf{憑證} & \textbf{用途} & \textbf{管理重點} \\
\hline
Channel Secret & 驗證 Webhook 簽章，確認事件內容未被竄改 & 在 Channel 的 Basic settings 檢視；若外洩，管理者重新簽發並同步更新伺服器 \\
\hline
Channel Access Token & 授權 Bot Server 呼叫 Messaging API & 從 Messaging API 設定取得或簽發；建議選擇符合需求的有效期限，疑似外洩時立即撤銷／重發 \\
\hline
\end{tabular}
\end{table}

\section{取得並設定憑證}

在 LINE Developers Console 選取 Messaging API Channel：Channel Secret 位於 Basic settings；Channel Access Token 位於 Messaging API 設定頁。可用的 Token 類型包含有效期限自訂、Stateless、短期及長效等，效期與簽發數量限制不同。官方文件建議使用者依部署需求選擇適當類型；不要為每次 API 呼叫重新簽發 Token，也不要將 Token 放在前端網頁、Git 儲存庫或公開日誌。

本範例從作業系統環境變數讀取兩項憑證。開發時可在 PowerShell 當前終端機設定暫存值：
\begin{lstlisting}[language=bash,caption={Windows PowerShell 設定本次工作階段的憑證環境變數},label={lst:env}]
$env:LINE_CHANNEL_SECRET = "your-channel-secret"
$env:LINE_CHANNEL_ACCESS_TOKEN = "your-channel-access-token"
\end{lstlisting}

上面的文字是佔位符，必須換成自己的 Channel 憑證，且只在可信任的本機終端機輸入。環境變數只作用於目前 PowerShell 工作階段及其子程序；關閉視窗後即消失。正式服務宜使用作業系統 Secret Store、雲端 Secret Manager 或部署平台的加密環境變數，不應把秘密寫入版本控制。

\chapter{Webhook 設定與測試}

\section{Webhook 的工作方式}

官方帳號好友傳送訊息或執行特定動作時，LINE 平台會向設定的 Webhook URL 發送 HTTPS POST。伺服器需讀取原始請求本文與 \texttt{X-Line-Signature}，使用 Channel Secret 驗證簽章；驗證成功後再解析事件。若事件是文字訊息，程式可用 \texttt{replyToken} 呼叫回覆 API，將訊息送回使用者。事件可能包含非文字型態，程式應明確忽略或各自處理。

\begin{figure}[htbp]
\centering
\renewcommand{\arraystretch}{1.45}
\begin{tabular}{c}
\fbox{\shortstack{使用者在 LINE 傳送訊息}}\\
$\downarrow$\\
\fbox{\shortstack{LINE 平台建立事件與簽章}}\\
$\downarrow$\\
\fbox{\shortstack{HTTPS POST 至 /callback}}\\
$\downarrow$\\
\fbox{\shortstack{以 Channel Secret 驗證原始本文}}\\
$\downarrow$\\
\fbox{\shortstack{解析文字事件並取出 replyToken}}\\
$\downarrow$\\
\fbox{\shortstack{呼叫 Reply API 回覆使用者}}
\end{tabular}
\caption{Webhook 接收、驗證與回覆的處理順序}
\label{fig:webhook-flow}
\end{figure}

\section{部署 HTTPS 端點}

Developers Console 的 Webhook URL 必須是可由 LINE 平台連線的 HTTPS 網址，並使用一般瀏覽器信任的 TLS 憑證。一般家用網路的 \texttt{localhost} 無法直接被 LINE 平台存取。開發測試可使用 ngrok 將本機 Flask 服務暫時轉送到 HTTPS 公開網址；正式環境則部署到具有效憑證的伺服器或雲端服務。

\begin{enumerate}
  \item 在本機安裝 Python 3.10 以上版本及 ngrok，並以 ngrok 帳號完成必要設定。
  \item 在一個終端機啟動稍後的 Flask 範例，監聽本機 5000 埠。
  \item 在另一個終端機執行 \texttt{ngrok http 5000}，複製 ngrok 顯示的 HTTPS Forwarding 網址。
  \item 在 Developers Console 的 Messaging API 分頁，編輯 Webhook URL，填入 \texttt{https://你的網域/callback} 並更新。
  \item 按下 Verify。確認程式可回應測試請求後，啟用 Use webhook。
  \item 掃描同一頁顯示的 QR Code 加官方帳號為好友，傳送文字確認 Bot 有回覆。
\end{enumerate}

ngrok 網址可能在重新啟動後改變，網址改變時需同步更新 Console。Webhook Verify 主要測試端點可達性，並不代表一般使用者事件處理、簽章驗證及回覆 API 都已通過，因此仍需以實際好友訊息測試。不要在生產環境沿用公開測試隧道。

\section{簽章驗證原則}

簽章使用 HMAC-SHA256 計算，簽署資料是收到的原始 Webhook request body。應使用 LINE 官方 SDK 的簽章驗證功能，不要先將 JSON 解析後重新序列化再驗證，因為空白、跳脫字元或欄位排列變化都會使簽章不符。缺少簽章或驗證失敗時，拒絕處理請求。LINE 不公開固定 Webhook IP 範圍，因此應驗證簽章，而不是依賴 IP 白名單。

\chapter{Python Flask 範例：接收及回覆文字訊息}

\section{安裝環境}

範例採用 Python 3.10 以上版本、Flask 與 LINE 官方 Python SDK v3。建議在專案資料夾建立虛擬環境，再於其中安裝相依套件：
\begin{lstlisting}[language=bash,caption={建立 Python 虛擬環境並安裝範例套件},label={lst:install}]
py -m venv .venv
.\.venv\Scripts\Activate.ps1
python -m pip install --upgrade pip
python -m pip install Flask line-bot-sdk
\end{lstlisting}

若 PowerShell 阻擋虛擬環境啟動指令，可依組織安全政策調整執行方式，或直接使用 \texttt{.venv\textbackslash Scripts\textbackslash python.exe} 執行程式。SDK 套件名稱為 \texttt{line-bot-sdk}，程式匯入路徑則使用 \texttt{linebot.v3}。

\section{完整範例程式}

程式\ref{lst:flask-bot} 建立 \texttt{/callback} POST 路由，由 SDK 驗證 Webhook 簽章與派送事件處理函式；只回覆文字訊息，其他型態略過。回覆內容採回聲方式，將使用者剛傳送的文字回傳。程式從環境變數讀取憑證，請勿以真實密鑰取代佔位值後提交程式碼。

\begin{lstlisting}[language=Python,caption={Flask LINE Bot：驗證 Webhook 並回聲回覆文字訊息},label={lst:flask-bot}]
import os

from flask import Flask, abort, request
from linebot.v3 import WebhookHandler
from linebot.v3.exceptions import InvalidSignatureError
from linebot.v3.messaging import (
    ApiClient,
    Configuration,
    MessagingApi,
    ReplyMessageRequest,
    TextMessage,
)
from linebot.v3.webhooks import MessageEvent, TextMessageContent

channel_secret = os.environ.get("LINE_CHANNEL_SECRET")
channel_access_token = os.environ.get("LINE_CHANNEL_ACCESS_TOKEN")
if not channel_secret or not channel_access_token:
    raise RuntimeError("Please set LINE_CHANNEL_SECRET and LINE_CHANNEL_ACCESS_TOKEN")

app = Flask(__name__)
handler = WebhookHandler(channel_secret)
configuration = Configuration(access_token=channel_access_token)


@app.post("/callback")
def callback():
    signature = request.headers.get("X-Line-Signature")
    if not signature:
        abort(400, description="Missing X-Line-Signature")

    # Keep the original body intact; the SDK verifies its signature before parsing.
    body = request.get_data(as_text=True)
    try:
        handler.handle(body, signature)
    except InvalidSignatureError:
        abort(400, description="Invalid webhook signature")

    return "OK"


@handler.add(MessageEvent, message=TextMessageContent)
def handle_text_message(event):
    with ApiClient(configuration) as api_client:
        messaging_api = MessagingApi(api_client)
        messaging_api.reply_message(
            ReplyMessageRequest(
                reply_token=event.reply_token,
                messages=[TextMessage(text=event.message.text)],
            )
        )


if __name__ == "__main__":
    app.run(host="0.0.0.0", port=5000, debug=False)
\end{lstlisting}

\section{程式重點說明}

\begin{enumerate}
  \item \textbf{讀取設定：}程式啟動時讀取兩項環境變數；若未設定就停止並顯示錯誤，避免服務在缺少憑證的狀態下誤導排查。
  \item \textbf{接收請求：}Flask 的 \texttt{/callback} 路由只接受 POST。\texttt{request.get\_data(as\_text=True)} 取得本文字串，傳給 SDK 前沒有解析或改寫 JSON。
  \item \textbf{驗證簽章：}\texttt{handler.handle(body, signature)} 先驗證 \texttt{X-Line-Signature}，驗證失敗會丟出 \texttt{InvalidSignatureError}，程式以 HTTP 400 拒絕。
  \item \textbf{比對事件型態：}裝飾器只將文字訊息事件交給 \texttt{handle\_text\_message}。圖片、貼圖、加入好友等其他事件不會進入此函式。
  \item \textbf{送出回覆：}以事件提供的 \texttt{reply\_token} 與 \texttt{TextMessage} 建立 Reply API 請求。Reply token 有時效且通常只能使用一次；應立即回覆，不能當成長期推播目標。
\end{enumerate}

\section{啟動與端對端測試}

先於 PowerShell 設定環境變數（程式\ref{lst:env}），再執行：
\begin{lstlisting}[language=bash,caption={啟動 Flask Bot 開發伺服器},label={lst:run}]
python main.py
\end{lstlisting}

若將程式另存為其他檔名，請將命令中的 \texttt{main.py} 改為實際檔名。接著另開終端機執行 \texttt{ngrok http 5000}，將 ngrok HTTPS 網址及 \texttt{/callback} 路徑填入 Channel 的 Webhook URL。Verify 成功並開啟 Use webhook 後，以 LINE 帳號加入官方帳號好友並傳送「你好」，應收到相同文字。範例在本文件中展示為完整程式；要實際運行，請複製程式碼為 \texttt{main.py}，依前述方式安裝套件並設定自己的憑證。

若收到訊息卻沒有回覆，檢查 Flask 終端機錯誤、Channel Access Token 是否有效、Webhook 是否啟用，以及 Manager 的自動回應設定是否與 Bot 衝突。若 Console Verify 失敗，檢查 HTTPS 網址是否正確、伺服器是否啟動、路徑是否包含 \texttt{/callback}，以及 ngrok 是否仍在運作。

\chapter{部署、安全與常見問題}

\section{安全與可靠性清單}

\begin{itemize}
  \item 使用環境變數或 Secret Manager 保存 Channel Secret 與 Access Token，限制只有服務程序及授權管理者可讀取。
  \item 將程式碼、日誌與截圖中的憑證遮蔽；若懷疑外洩，立即撤銷或重新簽發 Token，並更新執行環境。
  \item 對每個 Webhook 驗證簽章；只接受預期路由與方法，並採 HTTPS。
  \item 避免將完整個人訊息或 Token 寫入日誌。正式服務僅保存除錯及維運必要資料，並設置存取控制與保留期限。
  \item 設置合理逾時、錯誤處理、健康檢查及監控。LINE 平台可能重送未成功處理的事件；正式系統應考慮以事件識別碼去重，避免重複副作用。
  \item 不以大量訊息對 LINE 平台進行負載測試；遵守 API 速率限制、訊息額度與 LINE 開發規範。
\end{itemize}

\section{常見錯誤排查}

\begin{description}
  \item[簽章驗證失敗] 確認使用相同 Channel 的 Channel Secret，並以原始本文驗證；不要在驗證前格式化 JSON，檢查代理或中介層是否改寫本文。
  \item[401 或授權錯誤] 檢查 Access Token 是否剪貼完整、是否過期或已撤銷，以及是否屬於正確 Channel。
  \item[Webhook Verify 不成功] 確認 Flask 正在運行、URL 為可公開存取的 HTTPS、路徑正確，且防火牆或隧道未阻擋請求。
  \item[有收到事件但未回覆] 查看 SDK/API 錯誤與事件型態；reply token 可能已過期或使用過，並確認官方帳號的自動回應設定未造成混淆。
  \item[中文字或程式碼顯示異常] 確認終端機、原始檔皆使用 UTF-8；LaTeX 文件使用 XeLaTeX 並確認指定字型已安裝。
\end{description}

\section{費用、部署與後續擴充}

開發測試可從小量互動開始，但正式使用前需核對所在地區方案、每月訊息額度及計費條款。實際部署時可將 Flask 服務放在支援 HTTPS 的雲端平台或自有伺服器，設定正式網域、TLS 憑證、秘密管理、日誌與告警。若需要多事件並行或較複雜工作，宜快速驗證簽章後將工作放入佇列，並以事件 ID 進行去重；此架構需同時管理 reply token 時效及重試。

此回聲範例只示範基本文字互動。下一步可加入關鍵字路由、資料庫查詢、Flex Message、Rich Menu、使用者同意管理及多語系內容。擴充功能時仍需遵守最小權限、資料保護與平台政策。

\chapter{總結與操作檢核}

完成本教學的操作順序如下：
\begin{enumerate}
  \item 註冊 Business ID，建立 LINE Official Account。
  \item 從 Official Account Manager 啟用 Messaging API，並謹慎選擇 Provider。
  \item 在 Developers Console 確認 Channel，取得 Channel Secret 與 Access Token。
  \item 安裝 Python、Flask 與官方 LINE SDK，透過環境變數設定憑證。
  \item 啟動 Flask 和 HTTPS 隧道，設定 Webhook URL，按 Verify 並啟用 Use webhook。
  \item 加官方帳號為好友，傳送文字驗證 Bot 是否驗章並回覆。
\end{enumerate}

Messaging API 串接的關鍵在於清楚區分帳號、Provider、Channel 與後端伺服器的角色。只要正確完成帳號啟用、以安全方式管理憑證、提供可公開存取的 HTTPS Webhook，並在處理事件前驗證簽章，就能建立基本且可擴充的訊息服務。

\chapter{參考資料}

以下資料均為 LINE 官方文件或官方 SDK 專案；網頁內容與介面可能更新，查閱日期為 \today。
\begin{enumerate}
  \item\label{ref:getting-started} LINE Developers, “Get started with the Messaging API,”\\
  \url{https://developers.line.biz/en/docs/messaging-api/getting-started/}
  \item LINE Developers, “Messaging API overview,”\\
  \url{https://developers.line.biz/en/docs/messaging-api/overview/}
  \item LINE Developers, “Build a bot,”\\
  \url{https://developers.line.biz/en/docs/messaging-api/building-bot/}
  \item LINE Developers, “Verify webhook signature,”\\
  \url{https://developers.line.biz/en/docs/messaging-api/verify-webhook-signature/}
  \item LINE Developers, “Channel access token,”\\
  \url{https://developers.line.biz/en/docs/basics/channel-access-token/}
  \item LINE Developers, “Messaging API development guidelines,”\\
  \url{https://developers.line.biz/en/docs/messaging-api/development-guidelines/}
  \item LINE Developers, “Messaging API pricing,”\\
  \url{https://developers.line.biz/en/docs/messaging-api/pricing/}
  \item LINE Developers, “LINE Messaging API SDKs,”\\
  \url{https://developers.line.biz/en/docs/messaging-api/line-bot-sdk/}
  \item LINE Corporation, “LINE Messaging API SDK for Python,”\\
  \url{https://github.com/line/line-bot-sdk-python}
\end{enumerate}

\chapter{本文件 LaTeX 原始碼}

本章依作業要求列出本文件的完整 LaTeX 原始碼，包含文件設定、章節內容、表格、流程圖與程式碼清單。程式目錄中的本項目可用來快速查找原始碼。以下 \texttt{lstinputlisting} 僅將來源檔作為文字排版，不會再次執行其中的 LaTeX 命令。

\lstinputlisting[language={[LaTeX]TeX},caption={本文件完整 LaTeX 原始碼},label={lst:source},numbers=left,numberstyle=\tiny\color{codegray},basicstyle=\ttfamily\scriptsize,breaklines=true,breakatwhitespace=false,columns=fullflexible,frame=single,tabsize=2]{main.tex}

\end{document}
